% ===================================================================
% Evaluation section draft for the AppWeaver paper (IEEEtran).
% Grounded in: runs/20260508T130259_final (gpt-oss, 700 records) and
% runs/20260702T102407_787374Z (Claude Fable 5 pilot, planner cache off).
% Companion artifacts: results_table.tex, fig_workflow_accuracy.pdf,
% fig_quality_metrics.pdf, fig_latency.pdf, fig_model_comparison.pdf.
% ===================================================================

\subsection{Prototype and Setup}
We prototype AppWeaver on an OpenAirInterface (OAI) RAN connected to a
FlexRIC Near-RT RIC over E2, with monitoring and control exercised through
E2SM-KPM, E2SM-RC, and the slice, traffic-control (TC), MAC, RLC, PDCP, and
GTP service models. Each RIC application is wrapped by a domain agent that
publishes a machine-readable Agent Card at a standard discovery endpoint
(\texttt{/.well-known/agents/}) describing its skills, measurements, control
scope, resource requirements, and tool endpoints; MCP servers expose the
applications' monitoring and control functions as typed, callable tools. The
orchestration core is a LangGraph state machine that drives each intent
through an observe--diagnose--approve--act--verify--summarize loop, with the
agents coordinated over an A2A JSON-RPC bus. A scheduler mediates all control
actions through shared and exclusive resource leases (e.g., KPM subscriptions
are shared, while slice, TC, and RC control are exclusive), queueing
conflicting workflows instead of rejecting them, and an RC guardrail agent
gates enforced actions on platform health with optional human-in-the-loop
approval. All state, leases, and audit events persist in SQLite so workflows
and operator decisions remain traceable across restarts.

The planner LLM is pluggable. Our main batch uses \texttt{gpt-oss}, an open
model served locally via Ollama; a second configuration routes planning
through Claude (Fable~5) in headless mode. In both cases the LLM operates
strictly at the workflow-orchestration level: it selects and orders
applications but never touches E2 or A1 directly, and every generated plan
passes a policy-catalog validation layer that repairs or rejects references
to unknown applications or actions before execution.

\subsection{Performance Evaluation}
We evaluate workflow composition over a catalog of 50 operator intents (11
mapping to known workflow templates and 39 open-world intents with no
predefined template), spanning latency assurance, slice management,
transport QoS, and deliberately seeded cross-application conflicts. AppWeaver
(agent-card discovery and composition) is compared against six baselines:
deterministic rule routing (\emph{Manual rule}), launching a single
application (\emph{Single xApp}), unconstrained LLM planning
(\emph{LLM-only}), and three LangGraph variants with catalog repair
(human-approved \emph{LG-Override}, autonomous \emph{LG-Enforced}, and
\emph{LG-LLM} with LLM-driven recovery). Each intent runs twice per method,
yielding 700 records. Table~\ref{tab:appweaver-results} and
Figs.~\ref{fig:accuracy}--\ref{fig:latency} summarize the results.

Three findings stand out. First, agent-card discovery is the only approach
that generalizes: AppWeaver composes the correct workflow for 50\% of
open-world intents while retaining 100\% accuracy on known intents, whereas
every catalog-bound method collapses to 12.8\% open-world accuracy and
LLM-only planning reaches just 17.9\%. Since open-world intents dominate
realistic service requests, this lifts overall composition accuracy to 61\%
versus 30--32\% for the strongest baselines. Second, the governance layers do
not tax reliability: AppWeaver sustains 100\% tool-call success and policy
validity with zero hallucinated application references, and detects 96\% of
seeded cross-application conflicts (100\% for the catalog-repaired LangGraph
variants); its LLM-driven recovery resolved every failure it attempted
in-run. Third, adaptivity has a latency cost that remains compatible with
orchestration timescales: mean intent-to-execution latency is 3.2\,s for
AppWeaver, versus 47--315\,ms for catalog-bound LangGraph variants and
1\,ms for rule-based routing, i.e., well within the seconds-to-minutes
orchestration loop that AppWeaver targets while native E2/A1 loops continue
at their own timescales.

To assess sensitivity to the planner LLM, we replayed a pilot subset of the
catalog with Claude (Fable~5) as the planning model under stricter
conditions (planner caching disabled, so every request pays the full
reasoning cost). On the covered intent--method cells both models compose
identical, correct workflows with full tool-call success
(Fig.~\ref{fig:models}); end-to-end latency rises to 11--16\,s, reflecting
uncached remote-scale reasoning versus the locally served open model. This
suggests the Agent-Card abstraction, rather than a particular LLM, carries
the composition capability, though a full-catalog comparison across models
remains future work. Two further caveats apply: live before/after KPM
deltas were unavailable in this batch (the telemetry bus was offline), so we
report orchestration-level metrics rather than network-effect KPIs; and the
harness-level recovery rate includes external replay, so we cite only
in-run (internal) recovery for the recovery claim above.
